\documentclass[numbered,webpdf,modern,large]{oup-authoring-template}
\usepackage[mathlines,switch]{lineno}
\linenumbers
\usepackage{amsmath}
\usepackage{booktabs}
\usepackage{graphicx}
\usepackage{xurl}
\usepackage{xcolor}
\newcommand{\rev}[1]{\textcolor{red}{#1}}
\graphicspath{{figures/output/}}
\newcommand{\societylogo}{}

\begin{document}

\journaltitle{Bioinformatics}
\DOI{DOI added during production}
\copyrightyear{2026}
\pubyear{2026}
\vol{XX}
\issue{x}
\access{Published: date added during production}
\appnotes{Original Paper}
\firstpage{1}

\title[Layer Choice in Protein Language Models]{Depth Matters but the Final
Layer Suffices: \rev{Leakage-Free Evaluation of Layer Choice in Protein
Language Models for Mutation Effect Prediction}}

\author[1,2]{Yutao Guo}
\author[1,2]{Zihan Zhang}
\author[1,2]{Xuezhou Zhao}
\author[1,2]{Mengxi Chen}
\author[1,3,$\ast$]{Dan Wu}

\address[1]{\orgname{International College of Pharmaceutical Innovation, Soochow
University}, \orgaddress{Suzhou, \postcode{215123}, \country{China}}}
\address[2]{\orgname{School of Pharmacy and Biomolecular Sciences, RCSI}, \orgaddress{
Dublin, \postcode{D02 VN51}, \country{Ireland}}}
\address[3]{\orgname{Department of Chemistry, RCSI}, \orgaddress{Dublin,
Dublin, \postcode{D02 YN77}, \country{Ireland}}}

\corresp[$\ast$]{Corresponding author. \href{mailto:danwu@suda.edu.cn}{danwu@suda.edu.cn}}

\received{11 September 2026}{0}{2026}
\revised{19 September 2026}{0}{2026}

\abstract{\textbf{Motivation:} Protein language models such as ESM-2 encode
biochemical information across 33 transformer layers. Multi-layer embedding fusion is
often used under the assumption that different layers provide complementary
mutation-relevant information. However, the conditions under which this assumption
holds for mutation effect prediction remain poorly characterized.
\textbf{Results:} \rev{We benchmarked two 650M-parameter protein language models
(ESM-2 and the structure-aware SaProt) on the ProteinGym substitution benchmark
using per-layer probing under three feature arms and four train--test splits,
multi-scale pooling, logit-lens scoring, layer-wise similarity analysis, and a
position-lookup control. Under the standard random split, all 33 layers gave
nearly indistinguishable predictions ($\rho \approx 0.60$), but a control
regressor given only random position vectors reproduced this level
($\rho = 0.600$) and collapsed to $0.10$--$0.14$ under leakage-free splits
(paired drop $-0.46$, Wilcoxon $p = 5.2 \times 10^{-12}$): random-split layer
interchangeability is largely position memorization. Under the three
leakage-free splits (modulo, contiguous, position-grouped), performance rose
with depth for both models (mutation-specific embeddings up to
$\rho = 0.53$), the best layer lay in a narrow late band (layers 30--32), and
the best layer exceeded the final layer by at most $\Delta\rho = 0.009$;
across the 20 non-degenerate best-versus-final comparisons, 8 were equivalent
at the $\pm 0.01$ margin (bootstrap 90\% confidence interval inside the
margin), 12 were inconclusive, and none exceeded the margin; in the remaining
7 configurations the best layer was the final layer itself. A
single-substitution-only sensitivity analysis reproduces the depth trend
and the structure-token benefit, with the final layer remaining within
$0.01$ $\rho$ of the best layer for SaProt and within $0.02$ for ESM-2.
A logit-lens
analysis validated against the models' official logits pathway showed the
same depth trend zero-shot (final layer best: ESM-2 $\rho = 0.443$, SaProt
$\rho = 0.437$), consistent with published ProteinGym zero-shot levels. Real
Foldseek 3Di structure tokens improved SaProt under every split---slightly
under random ($+0.003$ to $+0.032$) and substantially under leakage-free
evaluation ($+0.06$ to $+0.12$, all $p \leq 10^{-4}$). No fusion strategy
improved on the final layer under leakage-free evaluation (paired differences
$\leq +0.018$ $\rho$, none significant; several significantly negative with
nonlinear readouts), whereas the same comparisons under the random split
showed small significant advantages for fusion ($+0.01$ to $+0.05$): the
apparent benefit of multi-scale pooling is protocol-dependent and does not
survive leakage-free evaluation.
Switching the readout mattered more than fusing layers, but only for ESM-2
(last-layer LightGBM $+0.045$ over Ridge, $p = 0.011$). Depth is important
for mutation effect
prediction, the final layer is already sufficient, and random-split
evaluation masks both facts.}
\textbf{Availability and Implementation:} All code and per-dataset results are
available at \url{https://github.com/schwarzeteeguo-creator/esm2-layer-invariance}.
\rev{An archived snapshot of the code and per-dataset results for this revision is
deposited at Zenodo (DOI: 10.5281/zenodo.22844290).} The benchmarking pipeline,
including masked-position feature extraction, multi-scale pooling, and logit lens
analysis, is implemented in Python using PyTorch and scikit-learn.}
\keywords{protein language models, layer selection, final-layer sufficiency,
mutation effect prediction, ProteinGym, position leakage, logit lens, CKA}

\maketitle

%=============================================================================
\section{Introduction}
%=============================================================================

Protein language models (PLMs) pre-trained via masked language modeling (MLM) on
millions of protein sequences have become indispensable for mutation effect
prediction~\cite{rives2021biological,meier2021language,lin2023evolutionary}.
Among these, ESM-2~\cite{lin2023evolutionary}
is one of the most widely adopted, with a 650M-parameter architecture comprising 33
transformer layers that produce 1280-dimensional per-residue hidden states. These
embeddings have been successfully applied to diverse tasks including variant effect
prediction~\cite{frazer2021disease,hie2021learning}, protein design, and functional
annotation~\cite{elnaggar2022prottrans,brandes2022proteinbert}, establishing ESM-2 as a
foundational tool in computational biology.

The internal organization of information across transformer layers has been
extensively studied in natural language processing, where the ``BERTology''
literature~\cite{rogers2021bertology,tenney2019bert} has revealed that different layers specialize
in distinct linguistic phenomena: lower layers encode surface-level syntax, middle
layers capture semantic content, and final layers refine task-specific
predictions~\rev{\cite{tenney2019bert,rogers2021bertology}}.
This hierarchical view has motivated analogous investigations in protein language
models. A growing body of work posits that different transformer layers in PLMs
similarly encode complementary biochemical information, and that combining
representations from multiple layers should therefore improve downstream tasks.
Kumar and Jha~\cite{kumar2025} demonstrated that mid-to-late ESM-2 layers (20--33)
outperform the final layer for kinase function prediction, while early layers
(6--14) have been shown to encode local physicochemical properties such as
hydrophobicity and secondary structure propensity~\cite{detlefsen2022}.
Alsamkary, Elshaffei et al.~\cite{elshaffei2025} proposed hierarchical pooling and
pooled attention architectures for fusing representations across multiple PLMs and
reported up to a 12\% Spearman improvement for protein--protein interaction
prediction. More broadly, the strategy of extracting and fusing embeddings from
multiple transformer layers or multiple PLMs has become common practice in
PLM-based pipelines~\cite{rao2019tape,ko2026fusion,si2024contact}, often under
the implicit assumption that each source contributes non-redundant
information.

These claims, however, rest on an assumption that has not been systematically tested:
that different ESM-2 layers genuinely encode \emph{complementary} information for
mutation effect prediction, rather than merely offering different views of the same
underlying signal. If layers are functionally similar under a given readout and
evaluation protocol, then multi-scale fusion may add cost without improving prediction.
If layer complementarity appears only under specific readouts or split protocols,
then layer selection should be treated as a conditional design choice rather than
a universal rule. Resolving this question carries practical
consequences: the computational cost of extracting and storing embeddings from
multiple layers is substantial for large-scale screening efforts, and the
proliferation of ad-hoc layer-selection heuristics complicates reproducibility.

To resolve this question, we conducted a large-scale benchmark on ProteinGym~\cite{notin2024proteingym},
a comprehensive collection of deep mutational scanning (DMS) experiments spanning
diverse protein families, assay types, and organismal contexts. We employed
three independent experimental paradigms (Figure~\ref{fig:concept}): (1) per-layer supervised probing via Ridge
regression, (2) multi-scale embedding pooling with controlled feature sets, and
(3) logit lens zero-shot scoring. In parallel, we quantified layer-wise representational
similarity through cosine and CKA analysis. Together, these approaches ask how
much mutation-relevant information each transformer layer carries, how that
information grows with depth, and how the answer depends on the evaluation
protocol. \rev{Our findings revise the premise of the multi-scale-fusion
literature on three points. First, the flat per-layer profiles produced by the
standard random split are largely position memorization: a lookup-table control
reproduces them entirely. Second, under leakage-free evaluation, predictive
information rises with depth, but the final layers are interchangeable among
themselves---the last layer is within $\Delta\rho \leq 0.01$ of the best layer
everywhere we measured---so depth matters, yet extracting the final layer alone
is sufficient. Third, no multi-scale fusion strategy improves on the final
layer under either evaluation protocol, whereas structure tokens contribute a
genuine, measurable benefit that only leakage-free evaluation can see.}

\begin{figure}[ht]
\centering
\includegraphics[width=\columnwidth]{figures/output/fig0_concept.pdf}
\caption{\textbf{Experimental design overview.}
The three complementary experimental paradigms employed in this study:
(\textbf{a}) Per-layer probing: Ridge regression on hidden states extracted from each
of the 33 transformer layers \rev{under three feature arms (masked wild-type,
unmasked wild-type, and mutant) and four train--test splits (random, modulo,
contiguous, and position-grouped), quantifying the predictive information content per
layer together with a random-position lookup control}.
(\textbf{b}) Multi-scale pooling: four embedding fusion strategies (Last Layer,
Mean L20--33, Concat 5L, Attention-weighted\rev{, with attention weights trained
within each cross-validation fold}) tested across four readout models
(Ridge, Random Forest, MLP, LightGBM) to assess whether combining layers improves
prediction.
(\textbf{c}) Logit lens: applying the final-layer LM head to intermediate hidden
states for zero-shot mutation scoring, directly measuring how token predictions
are refined across depth.
(\textbf{d}) Layer-wise cosine similarity analysis: quantifying representational
divergence independent of predictive performance.
Together, these paradigms provide converging evidence on whether different
transformer layers encode complementary or redundant information for mutation
effect prediction.}
\label{fig:concept}
\end{figure}

%=============================================================================
\section{Methods}
%=============================================================================

\subsection{Models and Data}

We used SaProt\_650M\_AF2~\cite{su2024saprot}, a structure-aware variant of ESM-2 650M
that extends the vocabulary to 446 structure-informed tokens (21 amino acids $\times$ 21
Foldseek 3Di states~\cite{vankempen2024foldseek}), as our primary model. For validation,
we also tested standard ESM-2 650M~\cite{lin2023evolutionary} in both the probing and
logit lens experiments, enabling direct comparison between structure-aware and
sequence-only representations. All experiments used ProteinGym's substitution DMS
benchmark~\cite{notin2024proteingym}. \rev{The reference release is ProteinGym
v1.3, which contains 217 substitution assays; our experiments use the 63-assay
LMDB snapshot that the SaProt repository distributes for supervised fitness
prediction (\url{https://github.com/westlake-repl/SaProt}). All 63 assays are
contained in the v1.3 benchmark, the subset was fixed before any experiment and
is identical across all arms, models, and splits, and the full assay list is
tabulated in Supplementary Table S1. Because each assay contributes one
train--test problem, all $n$ refer to assays: $n = 63$ for the random, modulo,
and position-grouped splits; six datasets whose mutated positions form a single
contiguous block have undefined Spearman correlations under the contiguous
split, which is therefore reported on the remaining 57.}
The pooling benchmark was run on a computationally tractable subset of
20 datasets because it crossed four pooling strategies with four readout models and
multiple cross-validation folds. We therefore treat the pooling analysis as a
controlled validation of the probing result rather than as a full ProteinGym-wide
estimate. \rev{The subset comprises the first 20 datasets in alphabetical order,
fixed identically for both models.} Per-dataset mutation counts range from under 100 to over 200,000.
\rev{Each DMS variant is exploded into one \emph{row} per substituted position
(multi-site variants are common in this benchmark: they are
84.6\% of variants and contribute 95.8\% of all analyzed rows;
Supplementary Table S1). Datasets were capped at 10,000 variants before row construction
(11 of 63 datasets) and at 5,000 rows inside the pipeline (26 of 63 datasets
reach this ceiling), both with a fixed seed shared across arms, splits, and
layers. For SaProt, we additionally prepared real Foldseek 3Di structure tokens by
assigning each wild-type residue the 3Di state of its aligned residue in the
AlphaFold Protein Structure Database entry for that sequence (structures available for
62 of 63 datasets; two aligned structures were truncated at the C terminus and the
missing residues were padded with the Foldseek gap state; the influenza nucleoprotein
dataset has no AFDB entry and was evaluated with structure-masked tokens). SaProt
results are reported for both structure-masked and real-3Di inputs
(per-dataset 3Di file lengths and gap counts in Supplementary Table
S5).}

\subsection{Experiment 1: Per-Layer Probing}

\rev{For each DMS dataset we extracted per-layer hidden states under three feature
arms. The \textbf{masked\_wt} arm masks each mutated position in its own forward
pass of the wild-type sequence (``M\#'' for SaProt, \texttt{<mask>} for ESM-2;
single-position masking), and each row receives the hidden state at that row's
substitution position---the conventional masked-position probing protocol.
The \textbf{unmasked\_wt} arm reads the wild-type token at the mutated position
without masking (one forward pass per dataset, gathered at all mutated
positions). The \textbf{mutant} arm feeds the full mutant sequence and reads the
hidden state at the mutated position, averaged over the variant's mutated
positions for multi-site variants; it is the only arm whose representation
differs between two mutations at the same position, and therefore the only one
that can carry mutation-specific information beyond position identity. Both
SaProt\_650M\_AF2 (structure-masked and real-3Di inputs) and standard ESM-2 650M
were evaluated under identical conditions.}

\rev{Each arm was evaluated under four train--test splits, all with five
cross-validation folds (four for the single four-position dataset under the
position-grouped split). The \textbf{random} split assigns whole variants to
folds at random (seeded), so all rows of one variant always share a fold. The
\textbf{modulo} split assigns each residue position to fold
$(\text{position} - 1) \bmod 5$ and the \textbf{contiguous} split divides the
sequence into five contiguous segments, both following the official ProteinGym
protocol; under these schemes a row's fold is determined by its own
substitution position, so the rows of a multi-site variant can fall into
different folds (1.10M of 1.30M multi-site variants straddle modulo folds;
0.36M straddle contiguous folds). These splits are leakage-free with respect to
positions but retain variant-level label sharing for multi-site variants; a
single-substitution-only sensitivity analysis
(Section~\ref{sec:single_site}) therefore re-runs the mutant arm restricted to
single-site variants. Finally, a \textbf{position-grouped} split (GroupKFold
over residue positions, so that all mutations at a given position share a fold)
serves as a strict control. For each layer $l \in \{0,\dots,32\}$ and
feature configuration (combined: the hidden state concatenated with 42 auxiliary
features, namely one-hot encodings of wild-type and mutant amino acids, the BLOSUM62
substitution score, and a normalized Grantham distance proxy; embedding\_only; and
auxiliary\_only), we trained Ridge regression with the regularization strength
selected per dataset by generalized cross-validation~\cite{golub1979generalized}
(grid $10^{-2}$ to $10^{3}$), standardization fit inside each training fold, and
reported Spearman rank correlation $\rho$. The single
dataset with fewer than five mutated positions (four positions) used four
position-grouped folds.}

\rev{To quantify how much probing performance is attributable to position identity
alone, we added a \textbf{random-position control}: each mutation is represented by a
fixed random vector indexed by its residue position (33 independent draws), probed with
the identical Ridge protocol with and without auxiliary features. This control contains
no sequence information; any performance it recovers measures position memorization.}

\subsection{Experiment 2: Multi-Scale Pooling}

\rev{We compared four embedding strategies using fixed auxiliary features,
evaluated on the \textbf{mutant} arm under the modulo split as the primary
leakage-free protocol and under the random split as a protocol contrast: (1)
\textbf{Last Layer}: the standard approach using only the
final transformer layer;
(2) \textbf{Mean L20--33}: mean pooling across layers 20--33, motivated by
Kumar and Jha~\cite{kumar2025};
(3) \textbf{Concat 5L}: concatenation of layers 6, 14, 20, 26, and 32 (0-indexed,
including the final layer), spanning the full representational depth; and
(4) \textbf{Attention-Weighted}: softmax-weighted pooling across all 33 layers, with
the weights trained within each cross-validation fold (120 Adam steps jointly
with a closed-form ridge readout on training rows only, seeded).}

\rev{For each strategy, we evaluated Ridge regression (regularized by generalized
cross-validation as above), Random Forest (50 estimators,
max depth 10), MLP (two hidden layers of 256 units, ReLU), and LightGBM (100
estimators, max depth 7). Every dataset and strategy used the same five outer
cross-validation folds as the probing experiment; standardization was fit inside
training folds, and the auxiliary
features were identical across strategies, isolating the effect of embedding choice.}

\subsection{Experiment 3: Logit Lens}

\rev{We applied the model's final-layer LM head to the hidden states of each
intermediate layer, producing pseudo-logits for zero-shot mutation scoring. All
computations ran in fp32, and for every dataset the final-layer pseudo-logits were
validated against the model's official logits pathway (maximum $|\Delta \log P| <
10^{-3}$; all 63 datasets passed for both models); no per-layer trend is
interpreted except through gate-passed results
(cf.\ the tuned lens literature~\cite{belrose2023eliciting}). For standard ESM-2, we
used the conventional log-ratio, $\log P(\text{mutant}) - \log P(\text{wild-type})$.
For SaProt with structure-masked input tokens ($\text{AA}\#$, where $\#$ denotes
masked structure), we marginalized over all 21 structure-token variants per amino
acid; with real 3Di inputs, the wild-type structure token at each position was
retained. We reported per-layer Spearman $\rho$ between these zero-shot scores and
experimental DMS fitness values across all 63 datasets.}

\subsection{Layer Similarity Analysis}

\rev{We computed the $33 \times 33$ cosine similarity matrix of per-layer hidden
states, averaged over 200 randomly selected residue positions across 10 datasets
(a 50-position variant is retained for continuity with the original analysis;
the larger sample reduces the known small-sample upward bias of CKA). Because
raw cosine is not invariant to rotations of the coordinate system, we
additionally computed linear CKA~\cite{kornblith2019similarity} and the
debiased (unbiased-HSIC) linear CKA
estimator~\cite{szekely2021debiased} on the same hidden-state matrices. The unit
of analysis for similarity-versus-distance statistics is the set of unique layer
pairs $(i, j)$ with $i < j$, diagonal excluded. A high CKA combined with a
near-zero cosine indicates a change of basis that preserves linearly decodable
information.}

\subsection{Statistical Evaluation}

We used three complementary procedures: (1) an \textbf{invariance score}, defined
as the standard deviation of per-layer mean Spearman $\rho$ (smaller indicates a
flatter layer profile);
\rev{(2) \textbf{equivalence testing}: for each model, arm, and leakage-free
split, we tested whether the mean per-dataset difference between the best layer
and the final layer falls within a margin of $\Delta\rho = 0.01$. The margin was
fixed before the equivalence analysis as the smallest difference that could
plausibly change a downstream layer-choice decision; it is a reporting
convention, not a task-derived value, and is applied identically everywhere.
The best layer was selected as the argmax of the layer-mean curve computed on
the same per-dataset results; this same-data selection biases the procedure
towards detecting a difference, which makes any equivalence verdict
conservative. A comparison was declared \emph{Equivalent} when the bootstrap
90\% confidence interval of the paired mean difference (10{,}000 resamples,
seeded percentile method) lay entirely within $\pm 0.01$, a \emph{Difference}
when it lay entirely outside, and \emph{Inconclusive} otherwise; two one-sided
tests (TOST) $p$-values and Cohen's $d_z$ are reported alongside. In 7 of the
27 model--arm--split combinations the best layer was the final layer itself
(no test defined); the remaining 20 comparisons constitute Supplementary Table
S4; and (3) \textbf{paired Wilcoxon signed-rank tests} for split contrasts
(random versus leakage-free), for real-3Di versus structure-masked SaProt
inputs, for each fusion strategy versus the last layer, and for readout
contrasts, each reported with bootstrap 90\% confidence intervals of the paired
mean difference.}

%=============================================================================
\section{Results}
%=============================================================================

\subsection{Representations Rotate While Linearly Decodable Information Is Preserved}

We first asked whether different transformer layers encode representations that are
merely different views of the same underlying information, or whether they capture
genuinely complementary features. To answer this, we computed the $33 \times 33$ layer-wise
cosine similarity matrix of per-layer hidden states, averaged over 50 randomly selected
residue positions across 10 datasets (Figure~\ref{fig:similarity}).

\begin{figure}[ht]
\centering
\includegraphics[width=0.95\columnwidth]{figures/output/fig2_similarity_heatmap.png}
\caption{\textbf{\rev{Layer-wise similarity matrices for SaProt (33$\times$33).}
Averaged over 10 datasets and 200 random residue positions per dataset.
(\textbf{a}) Cosine: adjacent layers are highly similar ($\cos=0.98$); mean similarity
decays monotonically with distance and collapses to $-0.01$ only between the most
distant layers (Layer~0 vs.~32). (\textbf{b}) Linear CKA over the same hidden states:
the most distant layers retain CKA 0.81, indicating a rotation that preserves linearly
decodable information rather than information loss.}}
\label{fig:similarity}
\end{figure}

The matrix revealed a smooth representational trajectory across layers. Adjacent layers
were highly similar \rev{(mean cosine 0.96--0.98)}, and \rev{cosine decayed
monotonically with layer distance, gently at first and steeply only at the
largest distances (Pearson $r$ on unique layer pairs: $-0.93$ for ESM-2,
$-0.50$ for SaProt). For SaProt the mean cosine by distance fell from 0.97
($d=2$) through 0.90 ($d=10$) and 0.81 ($d=20$) to 0.50 ($d=30$), collapsing to
$-0.008$ between the most distant layers (Layer~0 vs.\ Layer~32); ESM-2 decayed
more gently and remained positive throughout (0.93 at $d=2$, 0.58 at $d=31$,
0.54 at $d=32$). Crucially, near-zero raw cosine does not imply lost
information: the linear CKA between SaProt layers 0 and 32 was 0.78 (0.81 in
the 50-position variant) while their cosine was $-0.008$. The CKA--distance
profile is non-monotonic for SaProt: CKA dips to 0.67 around $d=20$ and rises
again to 0.78 at $d=32$, so the most distant layer pair retains more linearly
shared information than mid-depth pairs---consistent with residual streams
accumulating rather than discarding early-layer content. For ESM-2 the same
pair gives CKA 0.41 with cosine 0.54 (adjacent-layer CKA 0.95--0.97 for both
models). Increasing the sample from 50 to 200 positions lowered the ESM-2
distant-layer CKA from 0.55 to 0.41, confirming that small-sample CKA is
upward-biased at this dimensionality; the debiased unbiased-HSIC estimator is
unstable here ($n < d$; point estimates can exceed 1), so we rely on the
larger-sample standard estimator (both estimators at both sample sizes are
tabulated in Supplementary Table S9). The depth-wise transformation is therefore
well described as a rotation of the representation that preserves much
linearly decodable information (Figure~\ref{fig:probing}c).} This large
rotation of the representational geometry raises a sharp question: \emph{does this
representational transformation reflect complementary information encoded at different
depths, or is mutation-relevant predictive performance preserved under this transformation?}

To answer this, we conducted three converging experimental paradigms, each designed to
interrogate the information content of individual layers from a different angle.

\subsection{\rev{Random-Split Layer Interchangeability Is Largely Position Memorization}}

We first asked \rev{whether the near-flat random-split layer profile reflects
genuine representational redundancy. Under our pipeline the profile is
exactly this flat: for the masked\_wt
arm, per-layer mean $\rho$ spans
$0.603$--$0.606$ for both models with invariance scores $\sigma \approx 5 \times
10^{-4}$ to $1.2 \times 10^{-3}$ (Table~\ref{tab:probing}). A Ridge regressor given
\emph{random vectors indexed only by mutation position} plus the same auxiliary
features recovered $\rho = 0.600$ under the random split (Table~\ref{tab:control}),
matching the full masked-position PLM probing level. Under leakage-free splits the
same control collapsed to $\rho = 0.138$ (modulo), $0.096$ (contiguous, $n=57$), and
$0.135$ (position-grouped), a paired drop of $-0.46$ (random versus modulo, Wilcoxon
$p = 5.2 \times 10^{-12}$); the random vectors alone carried $0.557$ under the random
split and approximately zero elsewhere, and the auxiliary features alone fell from
$0.339$ to $0.19$--$0.23$. Because every mutation at the same position shares one
masked-position embedding, the random-split task is solvable by a position lookup
table; the flat layer profile is therefore largely an artifact of position
memorization rather than evidence that mutation information is distributed uniformly
across depth.}

\begin{table}[ht]
\centering
\caption{\rev{Random-position lookup control (63 datasets; contiguous $n = 57$).
Mean Spearman $\rho$ of Ridge regressors given no sequence information. The control
reproduces the random-split PLM probing level and collapses under leakage-free
splits.}}
\label{tab:control}
\begin{tabular}{lcccc}
\toprule
Features & Random & Modulo & Contiguous & Position-grouped \\
\midrule
Random position vector + aux & 0.600 & 0.138 & 0.096 & 0.135 \\
Random position vector only  & 0.557 & $-0.008$ & $-0.056$ & $-0.010$ \\
Auxiliary only               & 0.339 & 0.230 & 0.188 & 0.228 \\
\bottomrule
\end{tabular}
\end{table}

\rev{Table~\ref{tab:probing} and Figure~\ref{fig:probing}a report the per-layer
probing of all three arms,
both models, and all four splits ($n = 63$; contiguous $n = 57$). Two patterns are
stable across every leakage-free column. First, predictive information rises
with depth: every leakage-free column improves monotonically from layer 0 to
the late 20s, and the best layer sits in a narrow late
band---layers 30--32 are best in 26 of 27 model--arm--split combinations (layer
27 in the remaining one), while in the other 7 combinations the best layer is
the final layer itself. Second, the best layer is at most marginally better than
the final layer: across the 20 non-degenerate best-versus-final comparisons the
mean per-dataset difference never exceeded $\Delta\rho = 0.009$; with the
unified verdict rule (bootstrap 90\% confidence interval of the paired mean
difference entirely inside $\pm 0.01$), 8 comparisons were equivalent, 12 were
inconclusive, and none indicated a difference beyond the margin (Supplementary
Table S4). Because the best layer is selected on the same data, this procedure
is biased towards detecting a difference, making the equivalence verdicts
conservative. The
invariance score rose from $\sigma \approx 10^{-3}$ (masked\_wt, random) to
$4$--$7 \times 10^{-2}$ under leakage-free evaluation, yet remained well below the
cross-dataset standard deviation ($\approx 0.17$): performance varies far more
\emph{across} datasets than \emph{across layers}, but under leakage-free evaluation
the residual layer variation has a consistent direction, favoring the final layers.
The mutant arm exceeded the masked arms in every split (e.g., 0.681 versus 0.606
under random, 0.483 versus 0.459 under modulo for ESM-2), confirming that
mutation-specific signal exists and is recoverable once the representation actually
depends on the mutation.}

\begin{table}[ht]
\centering
\caption{\rev{Unified per-layer probing, combined features (mean Spearman
$\rho$; best layer in parentheses, followed by the final-layer value when the best
layer is not layer 32). $n = 63$; contiguous $n = 57$. Full 33-layer curves in
Supplementary Table S2; per-dataset results in Supplementary Table S3.}}
\label{tab:probing}
\resizebox{\columnwidth}{!}{%
\begin{tabular}{llcccc}
\toprule
Model & Arm & Random & Modulo & Contiguous & Pos.-grouped \\
\midrule
ESM-2 & masked\_wt   & 0.606 (32) & 0.459 (32) & 0.370 (32) & 0.450 (32) \\
ESM-2 & unmasked\_wt & 0.606 (32) & 0.470 (32) & 0.388 (32) & 0.461 (32) \\
ESM-2 & mutant       & 0.681 (32) & 0.483 (31) $\rightarrow$ 0.479 & 0.365 (31) $\rightarrow$ 0.360 & 0.477 (31) $\rightarrow$ 0.475 \\
\midrule
SaProt (mask) & masked\_wt   & 0.603 (31) $\rightarrow$ 0.603 & 0.391 (30) $\rightarrow$ 0.386 & 0.282 (31) $\rightarrow$ 0.274 & 0.385 (31) $\rightarrow$ 0.382 \\
SaProt (mask) & unmasked\_wt & 0.604 (32) & 0.412 (30) $\rightarrow$ 0.406 & 0.324 (31) $\rightarrow$ 0.320 & 0.403 (27) $\rightarrow$ 0.400 \\
SaProt (mask) & mutant       & 0.652 (23) $\rightarrow$ 0.651 & 0.465 (30) $\rightarrow$ 0.460 & 0.323 (32) & 0.461 (30) $\rightarrow$ 0.458 \\
\midrule
SaProt (3Di) & masked\_wt   & 0.606 (32) & 0.461 (31) $\rightarrow$ 0.460 & 0.394 (31) $\rightarrow$ 0.394 & 0.454 (31) $\rightarrow$ 0.453 \\
SaProt (3Di) & unmasked\_wt & 0.606 (27) $\rightarrow$ 0.605 & 0.473 (31) $\rightarrow$ 0.466 & 0.415 (31) $\rightarrow$ 0.406 & 0.467 (31) $\rightarrow$ 0.459 \\
SaProt (3Di) & mutant       & 0.684 (26) $\rightarrow$ 0.683 & 0.531 (31) $\rightarrow$ 0.525 & 0.451 (31) $\rightarrow$ 0.445 & 0.532 (31) $\rightarrow$ 0.524 \\
\bottomrule
\end{tabular}}
\end{table}

\subsection{\rev{Multi-Scale Pooling Only Appears Beneficial Under the Random Split}}

\rev{We ran the pooling benchmark (mutant arm, GCV-regularized Ridge,
fold-internal standardization, attention weights trained per fold) on the
20-dataset design subset for each model under the leakage-free modulo split
(Table~\ref{tab:pooling}; Figure~\ref{fig:probing}b) and, under the identical
pipeline, under the random split. Under leakage-free evaluation no fusion
strategy improved on the last layer. For SaProt (structure-masked inputs),
all four strategies occupied a narrow band for every readout (Ridge spread
across strategies 0.005; LightGBM 0.014). For ESM-2, spreads were likewise
small (Ridge 0.041; LightGBM 0.050), and the best single configuration was
the last layer with LightGBM ($\rho = 0.508$), exceeding every multi-scale
strategy under every readout. The per-dataset paired differences
(Table~\ref{tab:fusion}) confirm this: under the modulo split no
strategy--readout combination significantly beats the last layer, and with
LightGBM several combinations are significantly \emph{worse} (ESM-2 Mean
L20--33 $-0.050$; SaProt Concat 5L $-0.012$). Under the random split,
in contrast, the same comparisons show small but significant advantages for
fusion ($+0.009$ to $+0.050$; five of six strategy--readout combinations for
SaProt, Concat 5L for ESM-2 under both readouts). The apparent benefit of
multi-scale pooling is therefore protocol-dependent: it emerges when
positions leak between folds and disappears---occasionally reversing---under
leakage-free evaluation. Per-fold-trained attention pooling was never the
best strategy under the leakage-free protocol. The full grid of
paired differences across all splits and all four readouts (including
contiguous and the RF/MLP readouts) is Supplementary Table S6, and the
strategy means under all four splits are Supplementary Table S8.}

\begin{table}[ht]
\centering
\caption{\rev{Paired fusion-versus-last-layer differences (mutant arm, combined
features, 20 datasets per model). Cells give the mean per-dataset difference
$\Delta\rho$ (strategy $-$ last layer) with the bootstrap 90\% confidence
interval; \dag{} marks intervals excluding zero. Under the leakage-free
modulo split no strategy significantly beats the last layer, while the random
split produces small significant fusion advantages.}}
\label{tab:fusion}
\resizebox{\columnwidth}{!}{%
\begin{tabular}{lcccc}
\toprule
& \multicolumn{2}{c}{\textbf{ESM-2}} & \multicolumn{2}{c}{\textbf{SaProt}} \\
Strategy & Random & Modulo & Random & Modulo \\
\midrule
\multicolumn{5}{l}{\emph{Ridge readout}} \\
Mean L20--33 & $-0.001$ ($-0.004$, $+0.002$) & $+0.010$ ($-0.001$, $+0.023$) & $+0.003$ ($-0.001$, $+0.007$) & $+0.001$ ($-0.012$, $+0.014$) \\
Concat 5L    & $+0.012$ ($+0.005$, $+0.021$)\dag & $+0.018$ ($-0.008$, $+0.044$) & $+0.030$ ($+0.020$, $+0.039$)\dag & $+0.003$ ($-0.016$, $+0.021$) \\
Attention    & $-0.002$ ($-0.011$, $+0.007$) & $-0.024$ ($-0.053$, $+0.002$) & $+0.010$ ($+0.004$, $+0.016$)\dag & $-0.002$ ($-0.012$, $+0.007$) \\
\midrule
\multicolumn{5}{l}{\emph{LightGBM readout}} \\
Mean L20--33 & $-0.002$ ($-0.005$, $+0.002$) & $-0.050$ ($-0.065$, $-0.035$)\dag & $+0.030$ ($+0.019$, $+0.040$)\dag & $-0.009$ ($-0.018$, $+0.001$) \\
Concat 5L    & $+0.009$ ($+0.006$, $+0.012$)\dag & $-0.013$ ($-0.025$, $-0.002$)\dag & $+0.050$ ($+0.033$, $+0.067$)\dag & $-0.012$ ($-0.018$, $-0.006$)\dag \\
Attention    & $-0.002$ ($-0.010$, $+0.005$) & $-0.045$ ($-0.086$, $-0.012$)\dag & $+0.046$ ($+0.029$, $+0.063$)\dag & $-0.014$ ($-0.027$, $-0.002$)\dag \\
\bottomrule
\end{tabular}}
\end{table}

\begin{table}[ht]
\centering
\caption{\rev{Multi-scale pooling under the unified leakage-free pipeline (mutant
arm, modulo split, combined features; 20 datasets per model; mean $\rho \pm$ sd).
Spreads across strategies: ESM-2 Ridge 0.041, RF 0.057, MLP 0.045, LightGBM 0.050;
SaProt Ridge 0.005, RF 0.010, MLP 0.039, LightGBM 0.014.}}
\label{tab:pooling}
\begin{tabular}{lcccc}
\toprule
\textbf{ESM-2 650M} & Ridge & RF & MLP & LightGBM \\
\midrule
Last Layer (baseline) & $0.463 \pm 0.174$ & $0.476 \pm 0.173$ & $0.480 \pm 0.198$ & $\mathbf{0.508} \pm 0.189$ \\
Mean L20--33          & $0.473 \pm 0.180$ & $0.420 \pm 0.165$ & $0.436 \pm 0.179$ & $0.457 \pm 0.174$ \\
Concat 5L             & $0.481 \pm 0.183$ & $0.453 \pm 0.181$ & $0.435 \pm 0.170$ & $0.495 \pm 0.190$ \\
Attention-weighted    & $0.439 \pm 0.201$ & $0.419 \pm 0.209$ & $0.440 \pm 0.209$ & $0.463 \pm 0.208$ \\
\midrule
\textbf{SaProt\_650M\_AF2} & Ridge & RF & MLP & LightGBM \\
\midrule
Last Layer (baseline) & $0.432 \pm 0.163$ & $0.415 \pm 0.174$ & $0.445 \pm 0.168$ & $\mathbf{0.448} \pm 0.178$ \\
Mean L20--33          & $0.433 \pm 0.166$ & $0.415 \pm 0.166$ & $0.429 \pm 0.166$ & $0.439 \pm 0.182$ \\
Concat 5L             & $0.434 \pm 0.185$ & $0.407 \pm 0.173$ & $0.406 \pm 0.176$ & $0.436 \pm 0.183$ \\
Attention-weighted    & $0.429 \pm 0.164$ & $0.405 \pm 0.166$ & $0.426 \pm 0.164$ & $0.433 \pm 0.177$ \\
\bottomrule
\end{tabular}
\end{table}

\subsubsection{\rev{Readout Choice Interacts With Layer Choice Only for ESM-2}}

\rev{Under the leakage-free protocol, switching the readout on the last layer
produced the largest single improvement observed anywhere in the pooling
benchmark for ESM-2 (LightGBM over Ridge: paired $\Delta\rho = +0.045$,
90\% CI $+0.021$ to $+0.068$, Wilcoxon $p = 0.011$), whereas for SaProt the
same contrast was not significant ($+0.016$, 90\% CI $-0.016$ to $+0.048$,
$p = 0.57$; all readout contrasts in Supplementary Table S7). Pooling
comparisons that vary the feature construction,
scaling, regularization, or split protocol across cells can manufacture much
larger spurious contrasts than any of these effects; with the pipeline held
fixed, no strategy--readout configuration involving multi-scale features
exceeded the best single-layer configuration for either model.}

\begin{figure}[ht]
\centering
\includegraphics[width=\columnwidth]{figures/output/fig3_probing_pooling_summary.png}
\caption{\textbf{\rev{Per-layer probing, multi-scale pooling, and layer distance under a
single fixed protocol.}}
(\textbf{a}) \rev{33-layer probing (combined features): under the random split the
masked-arm profile is flat for both models ($\sigma \approx 10^{-3}$) and is reproduced by
the random-position lookup control (dotted, 0.600); under leakage-free splits (modulo,
mutant arm) performance is lower and the best layers are consistently late (layers 30--32,
$\Delta\rho \leq 0.009$ versus the final layer).}
(\textbf{b}) \rev{Pooling (mutant arm, modulo split, 20 datasets per model; four
strategies as dots, mean connected): ESM-2 spread across strategies 0.041 (Ridge) to 0.057
(RF); SaProt spread 0.005 (Ridge) to 0.039 (MLP); last-layer LightGBM is the best single
configuration; per-fold-trained attention pooling is never best.}
(\textbf{c}) \rev{Cosine similarity versus layer distance (unique layer pairs):
SaProt decays monotonically and reaches $-0.01$ only at the largest distances, while
ESM-2 remains positive; linear CKA between the most distant SaProt layers remains 0.81
(Figure~\ref{fig:similarity}), indicating rotation with preserved linear information
rather than loss.}}
\label{fig:probing}
\end{figure}

\subsection{Logit Lens \rev{Confirms Late-Layer Information Accumulation}}

\rev{The lens (fp32, per-dataset validation against the official
logits pathway, $n = 63$) produced monotonically improving zero-shot scores with
depth for both models (Figure~\ref{fig:logit_lens}). ESM-2 rose from $\rho = 0.031$
(layer 0) to $\rho = 0.443$ (layer 32), with only shallow fluctuations in early
layers; SaProt with real 3Di tokens rose from $-0.077$ (layer 0) through
approximately zero around layer 20 to $\rho = 0.437$ (layer 32). Both curves attain
their maximum at the final layer, agreeing with the leakage-free probing:
mutation-relevant information accumulates toward the final layers. Zero-shot
absolute levels remain
below supervised probing (0.44 versus up to 0.53), so the lens is best read as a
consistency check on layer-wise trends rather than a standalone benchmark.}

\subsection{\rev{Agreement With Published ProteinGym Results}}
\label{sec:anchors}

\rev{The ProteinGym v1.3 score release provides per-assay zero-shot
predictions for published models, including ESM-2 650M; 46 of our 63
assays have published scores, and Table~\ref{tab:anchors} places our
final-layer results beside the published values recomputed on exactly
those 46 assays. Our validated logit-lens final layer gives 0.438 (ESM-2)
and 0.440 (SaProt, real-3Di), inside the published 650M-class band on the
same assays (ESM-2 650M 0.472, GEMME 0.474, EVE 0.478, DeepSequence
0.457, MSA Transformer 0.447, ESM-1v 0.445, ESM-2 3B 0.448, ESM-2 150M
0.421), and our independent SaProt zero-shot reproduction scores 0.484;
the official release contains no per-assay SaProt column, so SaProt is
anchored by these two independent scorings of our own. The remaining gap
to the official ESM-2 value (0.438 versus 0.472) traces to scoring
protocol: the official scorer masks each mutated position in turn
(masked-marginal), whereas the lens scores the full mutant sequence with
all mutations in place, a difference that matters because 95.8\% of
analyzed rows belong to multi-site variants. For supervised anchoring,
the same release ships out-of-fold predictions for published supervised
methods under the official fold schemes our pipeline implements.
Recomputed on our assays: Kermut 0.605, ProteinNPT 0.566,
embedding-augmented MSA Transformer 0.572 under the modulo fold (0.45 to
0.56 across methods under contiguous). Our leakage-free probing (0.32 to
0.53) falls inside this published supervised band. The published
one-hot baseline without augmentation drops from 0.570 under the random
fold to $-0.003$ under the modulo fold, an external confirmation of the
position-memorization mechanism that our lookup control demonstrates
(Supplementary Table S10).}

\begin{table}[ht]
\centering
\caption{\rev{Published ProteinGym v1.3 results recomputed on the same assays
as this work (46 of 63 have published scores). Panel A: zero-shot, mean
per-assay Spearman $\rho$. Panel B: supervised, the release's official
out-of-fold predictions under the stated fold scheme. Our rows are
matched to the same 46 assays.}}
\label{tab:anchors}
\small
\begin{tabular}{lr}
\toprule
\multicolumn{2}{l}{\textbf{A. Zero-shot} ($n=46$)} \\
Method & $\rho$ \\
\midrule
ESM-2 650M (published) & 0.472 \\
GEMME & 0.474 \\
EVE (single) & 0.478 \\
DeepSequence (single) & 0.457 \\
MSA Transformer (single) & 0.447 \\
ESM-1v (single) & 0.445 \\
ESM-2 3B & 0.448 \\
ESM-2 150M & 0.421 \\
\midrule
Our lens: ESM-2 final layer & 0.438 \\
Our lens: SaProt (3Di) final layer & 0.440 \\
Our SaProt zero-shot reproduction & 0.484 \\
\bottomrule
\end{tabular}

\vspace{4pt}
\begin{tabular}{lrrr}
\toprule
\multicolumn{4}{l}{\textbf{B. Supervised} (official folds, $n=46$)} \\
Method & Random & Modulo & Contiguous \\
\midrule
Kermut & 0.711 & 0.605 & 0.561 \\
ProteinNPT & 0.697 & 0.566 & 0.525 \\
Emb.-aug.\ MSA Transformer & 0.641 & 0.572 & 0.536 \\
OHE, not augmented & 0.570 & $-0.003$ & $+0.027$ \\
\midrule
Our leakage-free probing (final layers) & -- & \multicolumn{2}{c}{0.32--0.53} \\
\bottomrule
\end{tabular}
\end{table}

\begin{figure}[ht]
\centering
\includegraphics[width=\columnwidth]{figures/output/fig1_logit_lens.png}
\caption{\textbf{Logit lens: zero-shot mutation scoring by layer ($n = 63$).}
\rev{Both models improve with depth and attain their maximum at the final layer:
ESM-2 (orange) rises from $\rho = 0.03$ (layer 0) to $\rho = 0.443$ (layer 32);
SaProt with real 3Di tokens (blue) rises from negative early-layer values to
$\rho = 0.437$ (layer 32). Shaded bands indicate $\pm 1$ standard deviation across
datasets. All computations in fp32 and validated against the official logits
pathway.}}
\label{fig:logit_lens}
\end{figure}

\subsection{Boundary Conditions}

\rev{The results above localize what the random split does and does not
measure. We now examine the two factors that moderate every conclusion:
evaluation protocol and input representation.}

\subsubsection{\rev{Leakage-Free Splits: What Is Lost and What Remains}}

\rev{The control and main-table results above already localize the leakage: the
random-position lookup table reproduces the random-split level ($0.600$ versus
$0.603$--$0.606$ for full probing) and loses it entirely under leakage-free splits
($0.10$--$0.14$), while the random vectors alone carry $0.557$ under the random
split and approximately zero elsewhere (Tables~\ref{tab:control}
and~\ref{tab:probing}; Figure~\ref{fig:position_level}). Among the leakage-free
splits, modulo and position-grouped give similar levels (e.g., SaProt real-3Di
mutant arm: $0.531$ and $0.532$) while contiguous is uniformly hardest ($0.451$),
consistent with sequence-contiguous segments correlating with local structural
context and mutational scope. Random mutation splitting, the de facto standard in DMS
benchmarking, therefore inflates absolute performance estimates by roughly
$0.15$--$0.20$ $\rho$ in these arms, masks both layer differences and
structure-token contributions, and can make multi-scale fusion appear
beneficial (Section~3.3). \rev{ProteinGym itself defines the modulo and
contiguous leakage-free splits used here, so reporting one of them alongside
the standard random split requires no new infrastructure---analogous to the
time-series community's practice of reporting both random and temporal
cross-validation~\cite{roberts2017crossvalidation}.}}

\begin{figure}[ht]
\centering
\includegraphics[width=\columnwidth]{figures/output/fig5_position_level.png}
\caption{\textbf{\rev{Random versus leakage-free split probing (SaProt, 63 datasets).}}
(\textbf{a}) \rev{Random mutation split: flat per-layer profile
($\sigma \approx 10^{-3}$), reproduced by the random-position lookup control
($\rho = 0.600$).} (\textbf{b})
\rev{Leakage-free splits (modulo, position-grouped, contiguous): performance
drops to $0.27$--$0.53$ depending on arm and model (Table~\ref{tab:probing};
minimum 0.274, SaProt structure-masked masked\_wt under contiguous), the
invariance score rises to
$4$--$7 \times 10^{-2}$, and the best layers shift consistently to layers 30--32.}}
\label{fig:position_level}
\end{figure}

\rev{Together, these boundary conditions refine the central findings: apparent
layer interchangeability under random splitting is dominated by position
memorization, leakage-free evaluation reveals a consistent depth trend whose
final layers are sufficient, the apparent benefit of multi-scale fusion is
protocol-dependent, and the benefit of structure tokens emerges mainly under
leakage-free evaluation.}

\subsection{\rev{Structure Tokens Help Mainly Under Leakage-Free Evaluation}}

\rev{Because SaProt can be run with real 3Di tokens or with structure-masked
tokens, the design allows a paired test, dataset by dataset and layer by
layer, of what structure tokens actually contribute (Table~\ref{tab:3di}). Under the
random split the two input modes were nearly identical (paired differences at layer
32: $+0.003$ masked\_wt, $+0.002$ unmasked\_wt, $+0.032$ mutant). Under every
leakage-free split the real-3Di inputs were consistently superior, with mean paired
advantages of $+0.059$ to $+0.122$ across arms and splits, all with Wilcoxon
$p \leq 10^{-4}$. Structure tokens therefore carry genuine mutation-relevant
information: the contribution is already positive under random splitting
($+0.003$ to $+0.032$, significant for the mutant arm) and becomes large when
position memorization is removed from the evaluation. This also
explains why structure-aware inputs add little in random-split comparisons:
the position channel, available to both input modes, saturates most of the
task.}

\begin{table}[ht]
\centering
\caption{\rev{Real-3Di versus structure-masked SaProt inputs (paired per-dataset
differences at layer 32, combined features; Wilcoxon signed-rank $p$). The benefit of
structure tokens is small and mostly non-significant under the random split and
large and highly significant under all leakage-free splits.}}
\label{tab:3di}
\begin{tabular}{lcc|ccc}
\toprule
& \multicolumn{2}{c|}{\textbf{Random}} & \multicolumn{3}{c}{\textbf{Leakage-free}} \\
Arm & $\Delta\rho$ & $p$ & Modulo & Contiguous & Pos.-grouped \\
\midrule
masked\_wt   & $+0.003$ & $2.2 \times 10^{-2}$ & $+0.073$ ($3.5 \times 10^{-7}$) & $+0.120$ ($8.7 \times 10^{-8}$) & $+0.071$ ($8.4 \times 10^{-7}$) \\
unmasked\_wt & $+0.002$ & $9.1 \times 10^{-1}$ & $+0.060$ ($4.6 \times 10^{-6}$) & $+0.086$ ($2.1 \times 10^{-6}$) & $+0.059$ ($1.2 \times 10^{-5}$) \\
mutant       & $+0.032$ & $2.5 \times 10^{-6}$ & $+0.064$ ($6.9 \times 10^{-5}$) & $+0.122$ ($6.5 \times 10^{-7}$) & $+0.067$ ($4.6 \times 10^{-5}$) \\
\bottomrule
\end{tabular}
\end{table}

\subsection{\rev{Sensitivity Analysis: Single-Substitution Variants Only}}
\label{sec:single_site}

\rev{Because 95.8\% of analyzed rows belong to multi-site variants whose
rows can share fitness labels across the official position-based folds
(Section~2.1), we re-ran the probing experiment for all three model
configurations restricted to single-substitution variants: one row per
variant, and no variant-level label sharing across folds remains. The
depth trend and the late-band optimum are unchanged---the best layer is
L30 or L31 in every leakage-free column (Supplementary Table S11).
Final-layer sufficiency is preserved for SaProt, where the
best-versus-final difference stays within $\Delta\rho \leq 0.010$
(structure-masked $\leq 0.006$, all bootstrap 90\% confidence intervals
overlapping zero; real-3Di $\leq 0.010$). For ESM-2 the gap widens
modestly: best-minus-final reaches $+0.016$ (modulo; 90\% CI $+0.006$ to
$+0.027$) and $+0.019$ (position-grouped; CI $+0.004$ to $+0.036$), so
on single-substitution data the ESM-2 final layer concedes up to about
$0.02$ $\rho$ to layers 30--31, slightly beyond the 0.01 full-benchmark
margin. The structure-token benefit replicates, and strengthens, on
single-site data: real-3Di over structure-masked at the final layer
gives $+0.028$ under random ($p = 4.3 \times 10^{-5}$) and $+0.057$ /
$+0.132$ / $+0.058$ under modulo / contiguous / position-grouped (all
$p \leq 10^{-3}$). Finally, an Envision-style supervised baseline
(one-hot and substitution features with MSA PSSM columns where
available, gradient-boosted readout; 50 of 63 assays with usable MSAs)
shows the same protocol signature externally: $\rho = 0.634$ under
random versus 0.410 modulo and 0.302 contiguous
(Supplementary Table S10).}

%=============================================================================
\section{Discussion}
%=============================================================================

\rev{The central result of this study has two halves. Under the standard
random-split protocol, the 33 transformer layers of SaProt and ESM-2 show
near-flat mutation effect prediction performance, and the position-lookup
control shows that this flatness is largely position memorization rather than
distributed mutation information. Under leakage-free evaluation the picture
inverts: mutation-specific performance rises with depth ($\rho$ up to 0.53),
the best layers are consistently the last few, the final layer is within
$0.01$ $\rho$ of the best layer in every full-benchmark comparison (within
$0.02$ on a single-substitution-only subset for ESM-2), and structure
tokens deliver a large measurable benefit ($+0.06$ to $+0.12$). Multi-scale
pooling provides no gain under leakage-free evaluation for either model; its
apparent benefit appears only under the leaking protocol. The practical
conclusion is not that multi-scale embeddings are never useful, but that for
mutation effect prediction with 650M PLMs the final layer is sufficient, and
that the evaluation protocol determines whether a benchmark can see the
difference.}

\subsection{\rev{Why Does Depth Improve Mutation-Effect Prediction?}}

\rev{The depth trend under leakage-free evaluation is consistent with the MLM
pre-training objective. All layers receive gradients through the same
cross-entropy loss computed at the final LM head, so each successive layer is
optimized to move representations closer to the quantities the LM head must
predict---conditional amino-acid distributions---which are what a
mutation-effect probe ultimately reads out. Layers therefore accumulate,
rather than specialize, mutation-relevant signal, and the last layers carry
the most of it. Consistent with accumulation rather than specialization, the
residual layer variation under leakage-free evaluation has a consistent
direction (later is better) instead of task-dependent optima at different
depths.}

\rev{The cosine similarity between layers decays monotonically with distance and,
for SaProt, collapses to approximately zero only between the most distant layers,
while the linear CKA between the same layers remains 0.78. This combination, a
near-rotation that preserves linearly decodable information, is consistent with the
effect of residual connections~\cite{he2016deep}: because residual streams accumulate
representations from all preceding layers, late layers retain a linearly accessible
trace of early-layer content even as their raw coordinates rotate.}

This interpretation is consistent with the logit lens
hypothesis~\cite{nostalgebraist2020,belrose2023eliciting},
which posits that intermediate layers in transformer language models progressively
refine token predictions toward the final output distribution, rather than
constructing qualitatively different representations. \rev{The lens curves
support this directly for both models: zero-shot scoring improves with depth
and peaks at the final layer.} The near-orthogonality of the most distant SaProt
layers ($\cos \approx -0.01$) \rev{alongside their high linear CKA (0.78)}
suggests that mutation-relevant signal
remains accessible while
representations move through successive coordinate frames.

Notably, this pattern contrasts with findings from NLP transformers, where
probing classifiers consistently recover different linguistic features from
different layer depths~\cite{tenney2019bert,rogers2021bertology}. In BERT, for example,
surface syntactic features peak in early layers while semantic features dominate in
middle-to-late layers. \rev{The single-peaked, monotonic profile of the PLM
curves may
reflect the fundamentally different nature of the learning signal in protein
versus natural language: whereas language contains hierarchical structure at
multiple timescales (morphology, syntax, semantics, discourse), the MLM task on
protein sequences operates on a single level of abstraction (amino acid identity),
potentially eliminating the pressure for mid-network specialization.}

\subsubsection{\rev{What Separates SaProt from Standard ESM-2?}}

\rev{Under the unified protocol the two models differ mainly in two measured
ways: how much their final representations retain early-layer geometry (ESM-2
distant-layer cosine 0.54 with CKA 0.41; SaProt cosine $-0.008$ with CKA
0.78), and the benefit conferred by structure tokens. The directly measured
difference is the conditional structure-token benefit: real 3Di inputs improve
leakage-free probing by $+0.06$ to $+0.12$ while leaving random-split performance
nearly unchanged. Structure information is genuinely present and genuinely
useful, but leakage-free evaluation is required to measure it.}

\rev{This framing carries a testable prediction: fine-tuning standard ESM-2
with structural objectives (e.g., contact prediction or fold classification)
should reproduce a similar leakage-free advantage over the sequence-only baseline.
This remains a mechanistic hypothesis rather than a direct result of the present
experiments.}

\subsection{Reconciling with Prior Layer-Selection Studies}

Our findings do not contradict the observation that different layers may be
optimal for different \emph{task types}. Kumar and Jha~\cite{kumar2025}
demonstrated that mid-to-late ESM-2 layers (20--33) outperform the final layer
for kinase function prediction, achieving 32\% higher clustering quality and
superior classification accuracy. This is compatible with our narrower conclusion:
different downstream tasks may induce different optimal readout directions, while
mutation effect prediction under random-split probing shows little layer dependence.
The key distinction is between \emph{task-level} specialization and
\emph{within-task} layer complementarity under a specified readout and evaluation
protocol.

Similarly, the 12\% Spearman improvement reported by Alsamkary, Elshaffei
et al.~\cite{elshaffei2025} for attention-based fusion architectures was
observed in the context of multi-chain protein--protein interaction prediction,
a task that differs substantially from per-residue mutation effect scoring.
The benefit of sophisticated pooling may therefore be task-dependent rather than
inherent to multi-scale representation fusion. For mutation effect prediction,
our results suggest that such improvements should be interpreted in light of
model type, readout choice, and split protocol: under leakage-free evaluation
no fusion strategy improved on the final layer for either model, and the
small fusion advantages that appear under random splitting track the
position-leakage regime rather than additional mutation information
(Section~3.3).

\subsection{Practical Implications}

For researchers using protein language model embeddings for mutation effect
prediction, our results support two practical recommendations that depend on
the model variant, readout model, and evaluation protocol:

\begin{enumerate}
    \item \textbf{Structure-aware models (SaProt, ESM-2 fine-tuned with structural
    objectives):} The final transformer layer is a strong default. \rev{In our
    benchmarks, multi-scale fusion provided no practical benefit over single-layer
    embeddings for either model under leakage-free evaluation}, and 1280 dimensions
    sufficed for both linear and nonlinear readouts.
    This reduces
    storage and computational costs by a factor of 5--10 compared to the common
    practice of concatenating 6--8 layers (up to 10,240 dimensions).

    \item \textbf{Standard PLMs (ESM-2, ProtBERT, etc.):} \rev{the final layer
    with a nonlinear readout is the strongest configuration we measured:
    last-layer LightGBM reached $\rho = 0.508$ ($+0.045$ over the same layer
    with Ridge, $p = 0.011$), the best pooling configuration overall. Layer
    concatenation offered no reliable gain under leakage-free evaluation
    ($+0.018$ with Ridge, not significant, and negative with LightGBM) at
    $4.9\times$ the feature dimension. For SaProt the readout contrast was not
    significant ($+0.016$, $p = 0.57$).}
\end{enumerate}

Table~\ref{tab:cost} quantifies the practical trade-offs. \rev{All pooling
strategies require one forward pass per variant with all 33 layers collected,
so embedding-extraction cost is identical across strategies and dominates the
end-to-end time ($\sim$3 minutes per dataset on an RTX 5090 for every
strategy). What Concat~5L inflates is the readout: its covariance matrix
grows from $1{,}322^2$ to $6{,}442^2$ entries (a $24\times$ penalty on the
Ridge solve) and its per-mutation storage grows $4.8\times$---in exchange for
no benefit under leakage-free evaluation on
either model, and a configuration (ESM-2 last-layer LightGBM) that exceeds it
on a single layer.} Storing Concat~5L embeddings for a typical DMS
benchmark ($\sim$10,000 variants) requires 258~MB versus 53~MB for a single
layer, a gap that widens linearly with dataset size and becomes substantial
for genome-scale studies.

\begin{table}[ht]
\centering
\caption{Computational cost comparison across pooling strategies.
``Per-mutation'' storage assumes float32 precision. Time estimates are
per-dataset averages on an NVIDIA RTX 5090, including embedding extraction
and readout training across 20 ProteinGym DMS subsets.
\rev{Best achievable $\rho$ from Table~\ref{tab:pooling} (mutant arm,
modulo split).}}
\label{tab:cost}
\resizebox{\columnwidth}{!}{%
\begin{tabular}{lcccc}
\toprule
& Last Layer & Mean L20--33 & Attention & Concat 5L \\
\midrule
Feature dimension                  & 1,322 & 1,322 & 1,322 & 6,442 \\
Storage per mutation (KB)          & 5.2   & 5.2   & 5.2   & 25.2 \\
Ridge design matrix (5,000 vars)   & $6.6 \times 10^6$ & same & same & $3.2 \times 10^7$ \\
Ridge covariance matrix entries    & $1.7 \times 10^6$ & same & same & $4.1 \times 10^7$ \\
Avg.\ time per dataset (min)       & $\sim$3 & $\sim$3 & $\sim$3 & $\sim$3 \\
\midrule
\multicolumn{5}{l}{\small \textbf{Best achievable $\rho$ (ESM-2 / SaProt):}} \\
\quad With Ridge                    & \rev{0.463 / 0.432} & \rev{0.473 / 0.433} & \rev{0.439 / 0.429} & \rev{0.481 / 0.434} \\
\quad With LightGBM                 & \textbf{\rev{0.508 / 0.448}} & \rev{0.457 / 0.439} & \rev{0.463 / 0.433} & \rev{0.495 / 0.436} \\
\bottomrule
\end{tabular}}
\end{table}

In both cases, computational resources may be better directed toward feature types
that are plausibly complementary to PLM embeddings:
\begin{enumerate}
    \item Structural features (e.g., AlphaFold-predicted pLDDT, PAE~\cite{jumper2021alphafold},
          or Foldseek 3Di states~\cite{vankempen2024foldseek}), which complement
          sequence-based embeddings;
    \item Sequence-level representations that capture global protein properties
          beyond per-residue features;
    \item Multiple sequence alignment (MSA)-based approaches~\cite{riesselman2018deepsequence,
          hopf2017mutation} that leverage evolutionary information beyond what PLMs encode.
\end{enumerate}

As discussed under Boundary Conditions (Section~3.5), \rev{the position-lookup
control shows that} standard evaluation protocols can conflate mutation-effect learning with
position-level memorization. We recommend that future DMS benchmarks adopt
position-aware evaluation as a standard complement to random splitting.

\subsection{Limitations}

Several limitations should be noted. First, we tested the 650M-parameter
ESM-2 variant and its structure-aware derivative SaProt. Smaller (35M, 150M)
and larger (3B, 15B) variants may show different degrees of layer specialization,
and this remains to be tested. Second, our
experiments are restricted to per-residue mutation effect prediction; for other
tasks such as kinase function prediction~\cite{kumar2025}, subcellular
localization, or protein--protein interaction scoring~\cite{elshaffei2025},
layer-wise differences and the benefits of multi-scale fusion may be more
pronounced. Third, our pooling benchmark used a 20-dataset validation subset,
which is smaller than the \rev{63}-dataset probing analysis. \rev{Fourth,
the SaProt pooling runs use structure-masked inputs, whereas the probing and logit
lens runs include real 3Di inputs, and one dataset (influenza nucleoprotein) lacks
an AFDB structure. Fifth, the official ProteinGym modulo and contiguous splits
are position-based, so the rows of a multi-site variant---95.8\% of all
analyzed rows---can share a fitness label across folds; these splits remove
position leakage but not variant-level label sharing, and our
single-substitution sensitivity analysis (Section~\ref{sec:single_site})
addresses exactly this caveat. Sixth, our}
feature ablation (embedding-only vs.\ auxiliary-only) indicates that learned embeddings
dominate predictive performance in this setting, but it does
not rule out the possibility that nonlinear combinations of auxiliary features
could provide stronger baselines.
\rev{Seventh,} the
maximum sample size of 10,000 \rev{variants} per dataset (subsampled to 5,000
\rev{rows} for regression\rev{; 11 of 63 datasets were variant-capped and 26
reach the row ceiling}) may limit sensitivity to
subtle layer-wise effects on the largest DMS screens ($>$200,000 mutations),
though the consistency of \rev{layer profiles} across dataset-size categories argues
against this concern. \rev{Eighth,} our linear probing approach can only detect
information that is linearly separable; nonlinear interactions between layers
could theoretically exist but would require a different experimental design to
identify. \rev{Ninth, small-sample CKA is upward-biased at our dimensionality
(50-position ESM-2 distant-layer CKA 0.55 versus 0.41 at 200 positions), and
the debiased unbiased-HSIC estimator is unstable when $n < d$ (point estimates
can exceed 1); we therefore report the larger-sample standard estimator.
Tenth,} position-level data splitting (Section~3.5) shows that
random mutation splits permit position-level leakage, \rev{which the lookup
control quantifies and which motivates our reporting of leakage-free splits
throughout. Finally, our equivalence margin}
($\Delta\rho = 0.01$) \rev{was fixed before the equivalence analysis as the
smallest difference that could plausibly change a downstream layer-choice
decision; it is a reporting convention, not a task-derived requirement, and it
is applied identically everywhere.} \rev{The best layer in each comparison is
selected on the same per-dataset results, which biases the procedure towards
detecting a difference and therefore makes the equivalence verdicts
conservative.}

\subsection{Conclusion}

\rev{For mutation effect prediction with 650M protein language models, depth
matters: under leakage-free evaluation, predictive information rises
monotonically through the network and the best layers are the last ones. But
the final layer is already sufficient---it is within $0.01$ $\rho$ of the best
layer in every full-benchmark comparison we ran (and within $0.02$ on a
single-substitution-only subset for ESM-2), and no multi-scale fusion
strategy improved
on it; the small fusion advantages that do appear are artifacts of the
position-leaking random split, the same regime in which a position lookup
table with no sequence information reproduces the apparent layer
interchangeability of the literature. Structure tokens deliver a genuine,
measurable benefit ($+0.06$ to $+0.12$) that only leakage-free evaluation can
see. Together these results replace the multi-scale-fusion heuristic for this
task with a simpler rule---extract the final layer, invest in the readout, and
evaluate under a position-aware split---and they argue for treating claims of
layer-specific representations in PLMs with
the same scrutiny that has become standard in the NLP
community~\cite{rogers2021bertology}.}

%=============================================================================
\section*{Acknowledgments}
We thank the International
College of Pharmaceutical Innovation, SUDA-RCSI for institutional support.

\section*{Funding}
\rev{This work was supported by the Jiangsu Tepin Foundation (Z241922) and
the Health Research Board (ILP-PHR-2024-014).}

\section*{Competing Interests}
\rev{The authors declare no competing interests.}

\section*{Data Availability}
Code and results are available at \url{https://github.com/schwarzeteeguo-creator/esm2-layer-invariance}.
\rev{An archived snapshot of the code and per-dataset results for this
revision is deposited at Zenodo (concept DOI
10.5281/zenodo.22844290; archived version v3.0.0, DOI 10.5281/zenodo.23264208).}
Raw ProteinGym data: \url{https://github.com/OATML-Markslab/ProteinGym}.

\newpage
\begin{thebibliography}{99}

% === References numbered by order of first citation in the text ===

\bibitem{rives2021biological}
Rives, A.\ \textit{et al.} (2021)
Biological structure and function emerge from scaling unsupervised learning to 250
million protein sequences.
\textit{Proceedings of the National Academy of Sciences}, \textbf{118}(15), e2016239118.

\bibitem{meier2021language}
Meier, J.\ \textit{et al.} (2021)
Language models enable zero-shot prediction of the effects of mutations on protein
function.
\textit{Advances in Neural Information Processing Systems}, \textbf{34}, 29287--29303.

\bibitem{lin2023evolutionary}
Lin, Z.\ \textit{et al.} (2023)
Evolutionary-scale prediction of atomic-level protein structure with a language model.
\textit{Science}, \textbf{379}(6637), 1123--1130.

\bibitem{frazer2021disease}
Frazer, J.\ \textit{et al.} (2021)
Disease variant prediction with deep generative models of evolutionary data.
\textit{Nature}, \textbf{599}(7883), 91--95.

\bibitem{hie2021learning}
Hie, B.\ \textit{et al.} (2021)
Learning the language of viral evolution and escape.
\textit{Science}, \textbf{371}(6526), 284--288.

\bibitem{elnaggar2022prottrans}
Elnaggar, A.\ \textit{et al.} (2022)
ProtTrans: Toward understanding the language of life through self-supervised learning.
\textit{IEEE Transactions on Pattern Analysis and Machine Intelligence}, \textbf{44}(10),
7112--7127.

\bibitem{brandes2022proteinbert}
Brandes, N.\ \textit{et al.} (2022)
ProteinBERT: A universal deep-learning model of protein sequence and function.
\textit{Bioinformatics}, \textbf{38}(8), 2102--2110.

\bibitem{rogers2021bertology}
Rogers, A., Kovaleva, O.\ and Rumshisky, A. (\rev{2020})
A primer in BERTology: What we know about how BERT works.
\textit{Transactions of the Association for Computational Linguistics}, \textbf{8},
842--866.

\bibitem{tenney2019bert}
Tenney, I., Das, D.\ and Pavlick, E. (2019)
BERT rediscovers the classical NLP pipeline.
\textit{Proceedings of the 57th Annual Meeting of the Association for Computational
Linguistics}, 4593--4601.

\bibitem{nostalgebraist2020}
nostalgebraist (2020)
Interpreting GPT: the logit lens.
\url{https://www.lesswrong.com/posts/AcKRB8wDpdaN6v6ru/}.

\bibitem{kumar2025}
Kumar, A.\ and Jha, I.P. (2025)
Layer probing improves kinase functional prediction with protein language models.
\textit{arXiv}, 2512.00376.

\bibitem{detlefsen2022}
Detlefsen, N.S.\ \textit{et al.} (2022)
Learning meaningful representations of protein sequences.
\textit{Nature Communications}, \textbf{13}(1), 1914.

\bibitem{elshaffei2025}
Alsamkary, H., Elshaffei, M.\ \textit{et al.} (2025)
Beyond simple concatenation: Fairly assessing PLM architectures for multi-chain
protein--protein interactions prediction.
\textit{arXiv}, 2505.20036.

\bibitem{rao2019tape}
Rao, R.\ \textit{et al.} (2019)
Evaluating protein transfer learning with TAPE.
\textit{Advances in Neural Information Processing Systems}, \textbf{32}.

\bibitem{notin2024proteingym}
Notin, P.\ \textit{et al.} (2024)
ProteinGym: Large-scale benchmarks for protein fitness prediction and design.
\textit{Advances in Neural Information Processing Systems}, \textbf{36}.

\bibitem{su2024saprot}
Su, J.\ \textit{et al.} (2024)
SaProt: Protein language modeling with structure-aware vocabulary.
\textit{International Conference on Learning Representations (ICLR)}.

\bibitem{vankempen2024foldseek}
van Kempen, M.\ \textit{et al.} (2024)
Fast and accurate protein structure search with Foldseek.
\textit{Nature Biotechnology}, \textbf{42}(2), 243--246.

\bibitem{roberts2017crossvalidation}
Roberts, D.R.\ \textit{et al.} (2017)
Cross-validation strategies for data with temporal, spatial, hierarchical, or
phylogenetic structure.
\textit{Ecography}, \textbf{40}(8), 913--929.

\bibitem{he2016deep}
He, K.\ \textit{et al.} (2016)
Deep residual learning for image recognition.
\textit{Proceedings of the IEEE Conference on Computer Vision and Pattern Recognition (CVPR)},
770--778.

\bibitem{jumper2021alphafold}
Jumper, J.\ \textit{et al.} (2021)
Highly accurate protein structure prediction with AlphaFold.
\textit{Nature}, \textbf{596}(7873), 583--589.

\bibitem{riesselman2018deepsequence}
Riesselman, A.J., Ingraham, J.B.\ and Marks, D.S. (2018)
Deep generative models of genetic variation capture the effects of mutations.
\textit{Nature Methods}, \textbf{15}(10), 816--822.

\bibitem{hopf2017mutation}
Hopf, T.A.\ \textit{et al.} (2017)
Mutation effects predicted from sequence co-variation.
\textit{Nature Biotechnology}, \textbf{35}(2), 128--135.

\bibitem{golub1979generalized}
\rev{Golub, G.H., Heath, M.\ and Wahba, G. (1979)
Generalized cross-validation as a method for choosing a good ridge parameter.
\textit{Technometrics}, \textbf{21}(2), 215--223.}

\bibitem{kornblith2019similarity}
\rev{Kornblith, S., Norouzi, M., Lee, H.\ and Hinton, G. (2019)
Similarity of neural network representations revisited.
\textit{Proceedings of the 36th International Conference on Machine Learning},
3519--3529.}

\bibitem{belrose2023eliciting}
\rev{Belrose, N. \textit{et al.} (2023)
Eliciting latent predictions from transformers with the tuned lens.
\textit{arXiv}, 2303.08112.}


\bibitem{szekely2021debiased}
\rev{Sz\'ekely, G.J., Wattenberg, M.\ and Sundararajan, M. (2021)
Feature space or sample space? A debiased kernel regression estimator.
\textit{Advances in Neural Information Processing Systems}, \textbf{34},
7154--7164.}

\bibitem{gray2018envision}
\rev{Gray, V.E., Hause, R.J., Luebeck, J., Shendure, J.\ and Fowler, D.M. (2018)
Quantitative missense variant effect prediction using large-scale mutagenesis data.
\textit{Cell Systems}, \textbf{7}(1), 136--144.}

\bibitem{ko2026fusion}
\rev{Ko, Y.S., Parkinson, J.\ and Wang, W. (2026)
Scalable embedding fusion with protein language models: insights from
benchmarking text-integrated representations.
\textit{Briefings in Bioinformatics}, \textbf{27}(1), bbag014.}

\bibitem{si2024contact}
\rev{Si, Y.\ and Yan, C. (2024)
Protein language model-embedded geometric graphs power inter-protein contact
prediction.
\textit{eLife}, \textbf{12}, RP92184.}

\end{thebibliography}

\end{document}
